% Generated from the list of items (see scratchpad generator). Each node records the Lean
% declaration(s) formalizing it. \notready marks placeholders (opaque objects or True-statements).

\section{Introduction and terminology}

\begin{definition}[Local field]
\label{def:local-field}
\lean{AnalyticTheory.IsLocalField}
A local field is a field $K$ with a non-trivial non-archimedean absolute value for which $K$ is complete.
\end{definition}

\begin{definition}[Function field]
\label{def:function-field}
\lean{AnalyticTheory.IsFunctionFieldOver}
$F\,|\,K$ is a function field: finitely generated of transcendence degree $1$ over $K$, with $K$ algebraically closed in $F$.
\end{definition}

\begin{definition}[Genus (placeholder)]
\label{def:genus}
\lean{AnalyticTheory.genus}
\uses{def:function-field}
\notready
The genus of a function field, defined via Riemann--Roch. Placeholder: not in Mathlib.
\end{definition}

\begin{definition}[Conservative]
\label{def:conservative}
\lean{AnalyticTheory.IsConservative}
\uses{def:genus}
$F\,|\,K$ is conservative if its genus does not change under finite algebraic constant field extensions.
\end{definition}

\begin{definition}[Elliptic function field]
\label{def:elliptic}
\lean{AnalyticTheory.IsEllipticOver}
\uses{def:conservative, def:function-field}
An elliptic function field is a conservative function field of genus $1$.
\end{definition}

\begin{definition}[Prime divisor of degree one (placeholder)]
\label{def:degree-one-prime}
\lean{AnalyticTheory.HasPrimeOfDegreeOne}
\notready
$F\,|\,K$ has a prime divisor of degree $1$. Placeholder predicate.
\end{definition}

\begin{definition}[Constant field extension]
\label{def:constant-extension}
\lean{AnalyticTheory.IsConstantExtension}
$FL$ is the constant field extension of $F$ by $L$: generated by the images of $F$ and $L$.
\end{definition}

\begin{definition}[Splitting field]
\label{def:splitting-field}
\lean{AnalyticTheory.IsSplittingField}
\uses{def:degree-one-prime, def:constant-extension}
$L$ is a splitting field of $F\,|\,K$ if $FL\,|\,L$ has a prime divisor of degree $1$.
\end{definition}

\begin{definition}[Elliptic subfield (placeholder)]
\label{def:elliptic-subfield}
\lean{AnalyticTheory.IsEllipticSubfield}
\uses{def:Fq}
\notready
An elliptic subfield of $F_K(q)$. Placeholder predicate.
\end{definition}

\section{Fields of periodic functions (\S 2)}

\begin{definition}[The field $\Omega$]
\label{def:omega}
\lean{AnalyticTheory.Omega}
$\Omega$ is an algebraic closure of $K$.
\end{definition}

\begin{definition}[Extension of the absolute value]
\label{def:omega-norm}
\lean{AnalyticTheory.omegaNorm, AnalyticTheory.omegaNormedField}
\uses{def:omega}
The absolute value of $K$ extends uniquely to $\Omega$ (spectral norm); $\Omega$ is normed through it.
\end{definition}

\begin{lemma}[Uniqueness of the extension]
\label{lem:ext-unique}
\lean{AnalyticTheory.omegaNorm_unique_extension}
\uses{def:omega-norm}
Any absolute value on $\Omega$ extending that of $K$ equals the spectral norm.
\end{lemma}

\begin{definition}[Convergent Laurent series]
\label{def:convergent-laurent}
\lean{AnalyticTheory.IsConvergentLaurent}
A Laurent series $\sum a_n X^n$ is convergent if $|a_n|r^n\to 0$ for every $r>0$.
\end{definition}

\begin{definition}[Holomorphic function over $K$]
\label{def:holomorphic}
\lean{AnalyticTheory.IsHolomorphicOver}
\uses{def:convergent-laurent}
$f:\Omega^\times\to\Omega$ is holomorphic over $K$ if it is given by a convergent Laurent series with coefficients in $K$.
\end{definition}

\begin{definition}[The ring $H_K$]
\label{def:H-K}
\lean{AnalyticTheory.holSubring}
\uses{def:holomorphic}
The holomorphic functions over $K$ form a ring $H_K$ (packaged as a subring of $\Omega\to\Omega$).
\end{definition}

\begin{lemma}[$H_K$ is a domain]
\label{lem:H-K-domain}
\lean{AnalyticTheory.holSubring_isDomain}
\uses{def:H-K}
$H_K$ is an integral domain.
\end{lemma}

\begin{definition}[The field $M_K$]
\label{def:M-K}
\lean{AnalyticTheory.MeroRing}
\uses{lem:H-K-domain}
$M_K$ is the quotient field of $H_K$ (meromorphic functions).
\end{definition}

\begin{definition}[$K$-structure on $H_K$ (placeholder)]
\label{def:const-algebra}
\lean{AnalyticTheory.holAlgebra}
\uses{def:H-K}
\notready
Constant functions give $K \hookrightarrow H_K$. Placeholder instance.
\end{definition}

\begin{definition}[Analytic valuation $w_\alpha$ (placeholder)]
\label{def:analytic-valuation}
\lean{AnalyticTheory.anVal}
\uses{def:M-K}
\notready
$w_\alpha$ is the order of vanishing at $\alpha$, with $w_\alpha(0)=+\infty$. Placeholder.
\end{definition}

\begin{definition}[Translation automorphism (placeholder)]
\label{def:translation}
\lean{AnalyticTheory.translation}
\uses{def:M-K}
\notready
$\tau_a:f(X)\mapsto f(a^{-1}X)$ on $M_K$. Placeholder.
\end{definition}

\begin{definition}[The field $F_K(q)$]
\label{def:Fq}
\lean{AnalyticTheory.FqSub, AnalyticTheory.Fq}
\uses{def:translation}
The $q$-periodic meromorphic functions form a subfield $F_K(q)$ of $M_K$.
\end{definition}

\begin{definition}[Divisors]
\label{def:divisor}
\lean{AnalyticTheory.Divisors}
\uses{def:analytic-valuation}
A divisor is an integer-valued function $m$ on $\Omega^\times$, finite in annuli and constant on $K$-conjugacy classes.
\end{definition}

\begin{definition}[Divisor of a function]
\label{def:divisor-of}
\lean{AnalyticTheory.divisorOf}
\uses{def:divisor, def:analytic-valuation}
The divisor $(f)$ of a meromorphic function $f$.
\end{definition}

\begin{definition}[$q$-periodic divisors]
\label{def:Dq}
\lean{AnalyticTheory.DivisorsPeriodic}
\uses{def:divisor}
$D_K(q)$: divisors with $m_{q\alpha}=m_\alpha$.
\end{definition}

\begin{definition}[$\bmod^\times q$]
\label{def:mod-times}
\lean{AnalyticTheory.ModTimes}
$a\equiv b\bmod^\times q$ means $b=aq^k$ for some $k\in\mathbb Z$.
\end{definition}

\begin{definition}[$K$-conjugacy mod $^\times q$]
\label{def:conj-mod-q}
\lean{AnalyticTheory.ConjModQ}
\uses{def:mod-times}
$\alpha,\beta$ are $K$-conjugate mod$^\times q$ if some conjugate of $\beta$ equals $q^k\alpha$.
\end{definition}

\begin{definition}[Degree $\deg_q$ (placeholder)]
\label{def:deg-q}
\lean{AnalyticTheory.degQ}
\uses{def:Dq}
\notready
The degree of a $q$-periodic divisor, formula (12). Placeholder; the formula is stated in \texttt{prop:deg-q-formula}.
\end{definition}

\begin{definition}[Jacobi image $\Phi_q$ (placeholder)]
\label{def:Phi-q}
\lean{AnalyticTheory.PhiQ}
\uses{def:Dq, def:mod-times}
\notready
The Jacobi image $\Phi_q(\mathfrak m)\in K^\times/q$, formula (12). Placeholder.
\end{definition}

\begin{definition}[Fundamental theta function]
\label{def:theta}
\lean{AnalyticTheory.theta}
\uses{def:Dq}
$\Theta(X)=\prod_{n\ge 0}(1-q^nX^{-1})\prod_{n<0}(1-q^{-n}X)$.
\end{definition}

\begin{definition}[Theta functions $\Theta_\alpha$]
\label{def:theta-alpha}
\lean{AnalyticTheory.thetaAlpha}
\uses{def:theta}
$\Theta_\alpha(X)=\Theta(\alpha^{-1}X)$.
\end{definition}

\begin{proposition}[Functional equation of $\Theta$]
\label{prop:theta-functional-eq}
\lean{AnalyticTheory.theta_functional_equation}
\uses{def:theta}
$\Theta(q^{-1}X)=-X\,\Theta(X)$ (formula (7)).
\end{proposition}

\begin{proposition}[Proposition 1 (Abel--Jacobi)]
\label{prop:abel-jacobi-1}
\lean{AnalyticTheory.abel_jacobi_one}
\uses{def:Phi-q, def:deg-q, def:Fq, def:divisor-of, prop:theta-functional-eq}
A $q$-periodic divisor is the divisor of a $q$-periodic function iff $\deg_q(\mathfrak m)=0$ and $\Phi_q(\mathfrak m)\equiv 1\bmod^\times q$; then the function is unique up to $K^\times$.
\end{proposition}

\begin{corollary}[Corollary to Proposition 1]
\label{cor:periodic-with-prescribed-zero}
\lean{AnalyticTheory.exists_periodic_function_with_valuation}
\uses{prop:abel-jacobi-1, def:conj-mod-q}
For every $\alpha\in\Omega^\times$ there is a $q$-periodic $f$ with $w_\alpha(f)=1$, and $w_\beta(f)=0$ for $\beta$ not conjugate to $\alpha$ mod$^\times q$.
\end{corollary}

\begin{proposition}[Formula (12)]
\label{prop:deg-q-formula}
\lean{AnalyticTheory.degQ_formula}
\uses{def:deg-q}
$\deg_q(\mathfrak m)=\sum[K(\alpha):K]m_\alpha$ and the Jacobi image, over representatives of conjugacy classes.
\end{proposition}

\begin{theorem}[Item I: $F_K(q)$ is a function field]
\label{thm:fq-function-field}
\lean{AnalyticTheory.fq_is_function_field}
\uses{def:Fq, prop:abel-jacobi-1}
$F_K(q)$ is a function field over $K$ with $K$ algebraically closed in it.
\end{theorem}

\begin{theorem}[Item II: valuations of $F_K(q)$]
\label{thm:fq-valuations}
\lean{AnalyticTheory.fq_valuation_eq_iff, AnalyticTheory.fq_valuation_surjective, AnalyticTheory.fq_valuations_classified}
\uses{thm:fq-function-field, def:analytic-valuation}
$w_\alpha=w_\beta$ on $F_K(q)$ iff $\alpha,\beta$ are $K$-conjugate mod$^\times q$; the value group is $\mathbb Z$; every valuation over $K$ is some $w_\alpha$.
\end{theorem}

\begin{theorem}[Item III: residue fields (placeholder)]
\label{thm:fq-residue}
\lean{AnalyticTheory.fq_residue_degree}
\uses{thm:fq-valuations}
\notready
The residue field of $w_\alpha$ is contained in $K(\alpha)$ and the prime has degree $[K(\alpha):K]$. Placeholder (statement recorded as $\mathrm{True}$).
\end{theorem}

\begin{theorem}[Item IV: $F_K(q)$ is elliptic]
\label{thm:fq-elliptic}
\lean{AnalyticTheory.fq_is_elliptic}
\uses{thm:fq-function-field, def:elliptic}
$F_K(q)$ is an elliptic function field over $K$.
\end{theorem}

\begin{definition}[Principal divisors $P_K(q)$]
\label{def:principal-divisors}
\lean{AnalyticTheory.principalDivisors}
\uses{def:divisor-of}
Divisors of $q$-periodic functions.
\end{definition}

\begin{definition}[Degree-zero divisors $D^0_K(q)$]
\label{def:D0q}
\lean{AnalyticTheory.divisorsDegZero}
\uses{def:Dq, def:deg-q}
$D^0_K(q)$: $q$-periodic divisors of degree $0$.
\end{definition}

\begin{definition}[Divisor class group $C^0_K(q)$]
\label{def:C0q}
\lean{AnalyticTheory.classGroupDegZero}
\uses{def:principal-divisors, def:D0q}
$C^0_K(q)=D^0_K(q)/P_K(q)$.
\end{definition}

\begin{theorem}[Item V: Abel--Jacobi isomorphism]
\label{thm:abel-jacobi-iso}
\lean{AnalyticTheory.abel_jacobi_iso}
\uses{def:C0q, def:Phi-q, thm:fq-elliptic}
The Jacobi map induces an isomorphism $C^0_K(q)\cong K^\times/q$.
\end{theorem}

\begin{definition}[Translation attached to a class (placeholder)]
\label{def:translation-class}
\lean{AnalyticTheory.translationOfClass}
\uses{def:C0q}
\notready
The algebraic translation automorphism attached to $c\in C^0_K(q)$. Placeholder.
\end{definition}

\begin{corollary}[Corollary Va]
\label{cor:translation-Va}
\lean{AnalyticTheory.translation_class_unique}
\uses{thm:abel-jacobi-iso}
For $a\in K^\times$ there is a unique class $c_a$ with $\Phi_q(c_a)=a\bmod^\times q$; $\tau_a(\mathfrak p)\sim\mathfrak p+c_a$ (the second part is not formalized).
\end{corollary}

\begin{corollary}[Corollary Vb]
\label{cor:translation-Vb}
\lean{AnalyticTheory.translation_analytic_eq_algebraic}
\uses{cor:translation-Va, def:translation-class}
The analytic and algebraic translation automorphisms coincide on $F_K(q)$.
\end{corollary}

\begin{remark}[Remark 1 (extension of constants)]
\label{rem:S2-ext-const}
\uses{def:Dq}
\notready
Embedding $D_K(q)\subset D_L(q)$ with ramification indices. Not formalized.
\end{remark}

\begin{remark}[Remark 2 (Galois constant extensions)]
\label{rem:S2-galois}
\uses{thm:fq-valuations}
\notready
Galois action on $M_L$, $D_L(q)$ and $\Phi_q$. Not formalized.
\end{remark}

\begin{remark}[Remark 3 (infinite constant extensions)]
\label{rem:S2-infinite}
\uses{rem:S2-galois}
\notready
Infinite algebraic constant extensions. Not formalized.
\end{remark}

\begin{remark}[Proposition 2 (missing from transcription)]
\label{rem:S2-prop2-missing}
\notready
The book refers to Proposition 2 ($L_K(q|\mathfrak m)$, Riemann--Roch spaces). It was not included in the LaTeX transcription and is not yet formalized.
\end{remark}

\begin{definition}[$\Theta$ in $M_K$ (placeholder)]
\label{def:theta-M}
\lean{AnalyticTheory.thetaM}
\uses{def:theta}
\notready
The function $\Theta$ viewed in $M_K$. Placeholder.
\end{definition}

\begin{definition}[The field $F_K(\bar q\mid X^m)$ (placeholder)]
\label{def:Fq-qbar-Xm}
\lean{AnalyticTheory.FqbarXm}
\uses{def:Fq, def:commensurable}
\notready
The image field of $\mu_m:\bar f(X)\mapsto\bar f(X^m)$. Placeholder.
\end{definition}

\section{The Weierstra\ss\ function and the absolute invariant (\S 3)}

\begin{definition}[Derivation $D$ (placeholder)]
\label{def:D-operator}
\lean{AnalyticTheory.Dop}
\uses{def:M-K}
\notready
$D=X\,\frac{d}{dX}$ on $M_K$. Placeholder.
\end{definition}

\begin{definition}[$\zeta=D\log\Theta$]
\label{def:zeta}
\lean{AnalyticTheory.zetaM}
\uses{def:D-operator, def:theta-M}
$\zeta(X)=D\log\Theta(X)$.
\end{definition}

\begin{definition}[The function $\wp$]
\label{def:wp}
\lean{AnalyticTheory.wpM}
\uses{def:zeta}
$\wp=-D\zeta$ is $q$-periodic with pole divisor $2\mathfrak p_1$.
\end{definition}

\begin{definition}[The function $\wp'$ (placeholder)]
\label{def:wp-prime}
\lean{AnalyticTheory.wpPrimeM, AnalyticTheory.wpPrime_spec}
\uses{def:wp, def:D-operator}
\notready
$\wp\'=\tfrac12(D\wp-\wp)$, defined in every characteristic (formula (18)).
\end{definition}

\begin{definition}[The series $s_k$]
\label{def:s-k}
\lean{AnalyticTheory.sK}
$s_k=\sum_{m>0}m^kq^m/(1-q^m)=\sum_n\big(\sum_{d\mid n}d^k\big)q^n$ (formula (35)).
\end{definition}

\begin{definition}[Coefficient $B$]
\label{def:coeffB}
\lean{AnalyticTheory.coeffB}
\uses{def:s-k}
$B=-5s_3$ (formula (37)).
\end{definition}

\begin{definition}[Coefficient $C$]
\label{def:coeffC}
\lean{AnalyticTheory.coeffC}
\uses{def:s-k}
$C=-\tfrac1{12}(5s_3+7s_5)$, an integral power series in $q$ (formula (38)).
\end{definition}

\begin{theorem}[Theorem VI]
\label{thm:VI}
\lean{AnalyticTheory.thmVI_generation, AnalyticTheory.thmVI_poles, AnalyticTheory.thmVI_relation, AnalyticTheory.thmVI_reflection}
\uses{def:wp, def:wp-prime, def:coeffB, def:coeffC, def:Fq, def:reflection}
$\wp,\wp\'$ generate $F_K(q)$; the poles are $2\mathfrak p_1$ and $3\mathfrak p_1$; they satisfy $\wp\'^2+\wp\wp\'=\wp^3+B\wp+C$; the reflection acts as stated.
\end{theorem}

\begin{definition}[Reflection (placeholder)]
\label{def:reflection}
\lean{AnalyticTheory.reflectionM}
\uses{def:M-K}
\notready
$\varepsilon:f(X)\mapsto f(X^{-1})$. Placeholder.
\end{definition}

\begin{remark}[Weierstraß form (char $\ne 2,3$)]
\label{rem:normal-form-3}
\lean{AnalyticTheory.thmVI_weierstrass_form}
\uses{thm:VI}
In characteristic $\neq 2,3$, $(D\wp)^2=4(\wp+\tfrac1{12})^3-g_2(\wp+\tfrac1{12})-g_3$.
\end{remark}

\begin{definition}[The polynomial $\Delta(B,C)$]
\label{def:Delta-q}
\lean{AnalyticTheory.deltaPoly}
$\Delta=B^2-C-64B^3+72BC-432C^2$ (formula (42)).
\end{definition}

\begin{definition}[Absolute invariant $j(q)$]
\label{def:jq}
\lean{AnalyticTheory.jInvQ}
\uses{def:Delta-q, def:coeffB, def:coeffC}
$j(q)=(1-48B)^3/\Delta$.
\end{definition}

\begin{lemma}[Lemma 1 (inversion)]
\label{lem:1}
\lean{AnalyticTheory.lemma1_inversion}
For $R$ with $|c_i|\le 1$, $u\mapsto 1/u+R(u)$ is a bijection from $0<|u|<1$ onto $|v|>1$ with $|1/u+R(u)|=1/|u|$.
\end{lemma}

\begin{theorem}[Theorem VII]
\label{thm:VII}
\lean{AnalyticTheory.thmVII_absinv, AnalyticTheory.thmVII_bijective}
\uses{def:jq, lem:1}
$j(q)=1/q+R(q)$ with $R$ universal, $|j(q)|=1/|q|>1$; every $j$ with $|j|>1$ equals $j(q)$ for a unique $q$.
\end{theorem}

\begin{remark}[Remark 2 (product formula)]
\label{rem:delta-product}
\lean{AnalyticTheory.remark3_delta_product}
\uses{def:Delta-q}
$\Delta=q\prod_{n>0}(1-q^n)^{24}$ in every characteristic (formula (45)).
\end{remark}

\begin{theorem}[Theorem VIII (Hasse invariant of $F_K(q)$)]
\label{thm:VIII}
\lean{AnalyticTheory.thmVIII_hasse}
\uses{thm:VII, def:hasse-invariant}
The Hasse invariant of $F_K(q)$ is trivial.
\end{theorem}

\begin{theorem}[Theorem VIIIa]
\label{thm:VIIIa}
\lean{AnalyticTheory.thmVIIIa_characterization}
\uses{thm:VII, thm:VIII, def:absolute-invariant, def:hasse-invariant, thm:VI}
An elliptic $F\,|\,K$ with a prime of degree $1$ is $K$-isomorphic to $F_K(q)$ for some $0<|q|<1$ iff $|j|>1$ and $\gamma$ is trivial.
\end{theorem}

\section{Semi-periodic functions (\S\S 4--5)}

\begin{definition}[Galois group $G$]
\label{def:galois-group}
\lean{AnalyticTheory.GalGroup}
The Galois group of $L\,|\,K$.
\end{definition}

\begin{definition}[Characters of $G$]
\label{def:characters}
\lean{AnalyticTheory.CharGroup}
\uses{def:galois-group}
Characters $\chi:G\to\mathbb{T}$ (values in the unit circle).
\end{definition}

\begin{definition}[The field $K_\chi$]
\label{def:K-chi}
\lean{AnalyticTheory.Kchi}
\uses{def:characters}
$K_\chi$ is the fixed field of $\ker\chi$.
\end{definition}

\begin{definition}[$q$-trivial characters]
\label{def:q-trivial}
\lean{AnalyticTheory.IsQTrivial}
\uses{def:K-chi}
$\chi$ is $q$-trivial if $q=N_{K_\chi|K}(Q)$ for some $Q\in K_\chi$.
\end{definition}

\begin{definition}[$H^1(G,L^\times/q)$ (placeholder)]
\label{def:H1-Lq}
\lean{AnalyticTheory.H1Lq, AnalyticTheory.eQ}
\uses{def:galois-group}
\notready
The cohomology group $H^1(G,L^\times/q)$ and the map $e_q$ of formula (51). Placeholder.
\end{definition}

\begin{lemma}[Lemma 2]
\label{lem:2}
\lean{AnalyticTheory.lemma2_H1_iso}
\uses{def:q-trivial, def:H1-Lq}
\notready
$e_q$ is injective with image the $q$-trivial characters: $H^1(G,L^\times/q)\cong\widehat G(q)$.
\end{lemma}

\begin{definition}[Character invariant (placeholder)]
\label{def:char-invariant}
\lean{AnalyticTheory.charInvariant}
\uses{def:characters, def:elliptic}
\notready
The character invariant $\chi_E$ of a split elliptic field. Placeholder.
\end{definition}

\begin{theorem}[Theorem IX]
\label{thm:IX}
\lean{AnalyticTheory.thmIX_realise, AnalyticTheory.thmIX_classification}
\uses{lem:2, def:char-invariant, def:splitting-field}
\notready
Elliptic fields of period $q$ split by $L$ correspond to pairs $\{\chi,\chi^{-1}\}$ of $q$-trivial characters.
\end{theorem}

\begin{theorem}[Theorem X (minimal splitting field)]
\label{thm:X}
\lean{AnalyticTheory.thmX_minimal_splitting}
\uses{thm:IX, def:K-chi}
$K_\chi$ is the unique minimal algebraic splitting field of $E$.
\end{theorem}

\begin{definition}[Semi-periodic triple $\mathbf Q=(L,\sigma,Q)$]
\label{def:semi-periodic-triple}
\lean{AnalyticTheory.SemiPeriodicTriple}
\uses{def:galois-group}
A generator $\sigma$ of $G$ and $Q\in L$ with $N(Q)=q$.
\end{definition}

\begin{definition}[The automorphism $\tau_Q\sigma$ (placeholder)]
\label{def:tau-Q-sigma}
\lean{AnalyticTheory.tauQSigma}
\uses{def:semi-periodic-triple, def:translation}
\notready
The operator $\tau_Q\sigma$ on $M_L$. Placeholder.
\end{definition}

\begin{definition}[The field $F_K(\mathbf Q)$]
\label{def:FKQ}
\lean{AnalyticTheory.FKQ}
\uses{def:tau-Q-sigma}
Fixed field of $\tau_Q\sigma$ in $M_L$ (semi-periodic functions).
\end{definition}

\begin{theorem}[Theorem XI]
\label{thm:XI}
\lean{AnalyticTheory.thmXI_semi_periodic}
\uses{def:FKQ, thm:IX}
\notready
$F_K(\mathbf Q)$ is elliptic of period $q$, split by $L$, with character invariant $\chi(\sigma)=\exp(2\pi i/n)$. Statement recorded as $\mathrm{True}$.
\end{theorem}

\begin{theorem}[Theorem XII]
\label{thm:XII}
\lean{AnalyticTheory.thmXII_elliptic}
\uses{def:FKQ}
\notready
$F_K(\mathbf Q)$ is elliptic over $K$, with constant extension $F_L(q)$ over $L$.
\end{theorem}

\begin{theorem}[Theorem XIII]
\label{thm:XIII}
\lean{AnalyticTheory.thmXIII_prime_degrees}
\uses{thm:XII}
\notready
Prime divisors of $F_K(\mathbf Q)$ correspond to $L$-conjugacy classes mod$^\times Q\sigma$, with residue field $L(\alpha)$.
\end{theorem}

\begin{theorem}[Theorem XIV]
\label{thm:XIV}
\lean{AnalyticTheory.thmXIV_divisors}
\uses{thm:XIII}
\notready
The divisor group of $F_K(\mathbf Q)$ is $D_K(\mathbf Q)$; each divisor is an $N_{\mathbf Q}$-norm.
\end{theorem}

\begin{corollary}[Corollary XIIIa (index)]
\label{cor:XIIIa}
\lean{AnalyticTheory.corXIIIa_index}
\uses{thm:XIII}
\notready
$F_K(\mathbf Q)$ has index $n$.
\end{corollary}

\begin{corollary}[Corollary XIIIb (minimal splitting field)]
\label{cor:XIIIb}
\lean{AnalyticTheory.corXIIIb_minimal}
\uses{thm:XIII}
\notready
$L$ is the minimal algebraic splitting field of $F_K(\mathbf Q)$.
\end{corollary}

\begin{theorem}[Theorem XV (Abel--Jacobi for semi-periodic functions)]
\label{thm:XV}
\lean{AnalyticTheory.thmXV_abel_jacobi}
\uses{thm:XIV, prop:abel-jacobi-1}
\notready
$C^0_K(\mathbf Q)\cong N L^\times/q$.
\end{theorem}

\begin{corollary}[Corollary XVa]
\label{cor:XVa}
\lean{AnalyticTheory.corXVa_quotient}
\uses{thm:XV}
\notready
$C^0_K(\mathbf Q)\subset C^0_K(q)$ with quotient $K^\times/NL^\times$.
\end{corollary}

\begin{theorem}[Theorem XVI (norm theorem)]
\label{thm:XVI}
\lean{AnalyticTheory.thmXVI_norm}
\uses{thm:XIV, thm:XV}
\notready
Every semi-periodic $f\neq 0$ is the norm of some $g\in F_L(q)$.
\end{theorem}

\begin{theorem}[Table of cohomology groups]
\label{tab:cohomology}
\lean{AnalyticTheory.table_cohomology}
\uses{thm:XVI}
\notready
$H^0$ and $H^1$ of $D_L(q)$, $D^0_L(q)$, $C^0_L(q)$, $P_L(q)$, $F_L(q)^\times$. Recorded as $\mathrm{True}$.
\end{theorem}

\begin{corollary}[Corollary XVIa (Brauer)]
\label{cor:XVIa}
\lean{AnalyticTheory.corXVIa_brauer}
\uses{thm:XVI}
\notready
$\mathrm{Br}(F_K(\mathbf Q))\to\mathrm{Br}(F_L(q))$ is injective.
\end{corollary}

\begin{corollary}[Corollary XVIb]
\label{cor:XVIb}
\lean{AnalyticTheory.corXVIb_splitting}
\uses{cor:XVIa}
\notready
An algebra is split by $F_K(\mathbf Q)$ iff it is split by $L$.
\end{corollary}

\begin{theorem}[Theorem XVII]
\label{thm:XVII}
\lean{AnalyticTheory.thmXVII_degree}
\uses{lem:XVII}
\notready
Degrees, inseparability, automorphisms and Galois criterion of $F_K(q\mid X)$ over $F_K(\bar q\mid X^m)$. Recorded as $\mathrm{True}$.
\end{theorem}

\begin{theorem}[Theorem XVIII]
\label{thm:XVIII}
\lean{AnalyticTheory.thmXVIII_subfields}
\uses{thm:XVII, def:elliptic-subfield}
Every elliptic subfield of $F_K(q)$ is $F_K(\bar q\mid X^m)$ for a unique period $\bar q$ and integer $m$.
\end{theorem}

\begin{corollary}[Corollary XVIIIa (isogenies)]
\label{cor:XVIIIa}
\lean{AnalyticTheory.corXVIIIa_isogeny}
\uses{thm:XVIII, def:isogeny}
$F_K(q)$ receives an isogeny from $E$ iff $E\cong F_K(\bar q)$ with $\bar q$ commensurable with $q$.
\end{corollary}

\begin{corollary}[Corollary XVIIIb (no CM)]
\label{cor:XVIIIb}
\lean{AnalyticTheory.corXVIIIb_no_cm}
\uses{def:complex-multiplication}
$F_K(q)$ does not have complex multiplication.
\end{corollary}

\begin{theorem}[Theorem XIX]
\label{thm:XIX}
\lean{AnalyticTheory.thmXIX_subfields}
\uses{thm:XVIII, thm:IX, thm:X}
\notready
Elliptic subfields of $F$ with period $q$ correspond to integers $m$ with $[F:E]=dm$ and $\bar\chi=\chi^d$. Recorded as $\mathrm{True}$.
\end{theorem}

\section{Elliptic subfields and isogenies (\S 6)}

\begin{definition}[Commensurable periods]
\label{def:commensurable}
\lean{AnalyticTheory.Commensurable}
$\bar q$ is commensurable with $q$ if $\bar q^{\,d}=q^m$ for nonzero integers $d,m$ (formula (71)).
\end{definition}

\begin{lemma}[Inclusion (72) and symmetry (73)]
\label{lem:XVII}
\lean{AnalyticTheory.fq_qbar_Xm_subset}
\uses{def:Fq-qbar-Xm, def:commensurable}
$F_K(\bar q\mid X^m)\subseteq F_K(q\mid X)$, and $F_K(\bar q\mid X^{-m})=F_K(\bar q\mid X^m)$.
\end{lemma}

\begin{definition}[Isogeny]
\label{def:isogeny}
\lean{AnalyticTheory.IsIsogeny}
An isogeny is a $K$-embedding of elliptic function fields.
\end{definition}

\begin{definition}[Complex multiplication (placeholder)]
\label{def:complex-multiplication}
\lean{AnalyticTheory.HasComplexMultiplication}
\notready
$M(F)\neq\mathbb Z$. Placeholder predicate.
\end{definition}

\begin{proposition}[Absolute integrality of $j$]
\label{prop:abs-integral}
\lean{AnalyticTheory.prop_cm_absolutely_integral}
\uses{def:complex-multiplication, def:absolute-invariant-gen}
Over any field, complex multiplication forces $j$ to be absolutely integral.
\end{proposition}

\begin{corollary}[Corollary XIXa]
\label{cor:XIXa}
\lean{AnalyticTheory.corXIXa_galois}
\uses{thm:XIX}
\notready
Inseparability degree $p^r$, automorphisms $U/q$, and Galois criterion for $F|E$.
\end{corollary}

\begin{corollary}[Corollary XIXb]
\label{cor:XIXb}
\lean{AnalyticTheory.corXIXb_maximal_split}
\uses{thm:XIX}
\notready
A unique maximal elliptic subfield with period $\bar q$ splitting over $K$.
\end{corollary}

\begin{corollary}[Corollary XIXc]
\label{cor:XIXc}
\lean{AnalyticTheory.corXIXc_maximal_jacobian}
\uses{cor:XIXb}
\notready
A unique maximal subfield $F_0\subseteq F$ isomorphic to $F_K(q)$, with $[F:F_0]=n^2$.
\end{corollary}

\begin{corollary}[Corollary XIXd]
\label{cor:XIXd}
\lean{AnalyticTheory.corXIXd_degree_n}
\uses{cor:XIXb}
\notready
Degrees divisible by $n$; equality $[F:E_0]=n$ iff $q$ is an $n$-th power in $K$.
\end{corollary}

\begin{proposition}[Absolute invariant in Weierstraß form (A1.8)]
\label{prop:absinv-weierstrass}
\lean{AnalyticTheory.absInv_weierstrass}
\uses{def:absolute-invariant, def:normal-forms}
In Weierstraß normal form, $j=12^3g_2^3/(g_2^3-27g_3^2)$.
\end{proposition}

\section{Appendix \S A1: fields with a prime divisor of degree 1}

\begin{definition}[Equivalence modulo squares]
\label{def:square-equiv}
\lean{AnalyticTheory.SquareEquiv}
$a\sim b$ iff $a=t^2b$ for some $t\in K$.
\end{definition}

\begin{definition}[Equivalence modulo $\pi(K^+)$]
\label{def:pi-equiv}
\lean{AnalyticTheory.PiEquiv}
In characteristic 2: $a\sim b$ iff $a=z^2+z+b$.
\end{definition}

\begin{definition}[Normal forms (A1.3), (A1.6), (A1.26), (A1.32)]
\label{def:normal-forms}
\lean{AnalyticTheory.IsNormalFormA3, AnalyticTheory.IsWeierstrassNormal, AnalyticTheory.IsNormalFormA1, AnalyticTheory.IsNormalFormA32}
The normal forms used in §A1.
\end{definition}

\begin{definition}[Absolute invariant (placeholder)]
\label{def:absolute-invariant}
\lean{AnalyticTheory.absInv}
\uses{def:normal-forms}
\notready
The absolute invariant $j$ of $F\,|\,K$ with a prime of degree $1$. Placeholder.
\end{definition}

\begin{definition}[Hasse invariant (placeholder)]
\label{def:hasse-invariant}
\lean{AnalyticTheory.hasseInv}
\uses{def:normal-forms, def:square-equiv, def:pi-equiv}
\notready
The Hasse invariant $\gamma$ (representative). Placeholder.
\end{definition}

\begin{proposition}[Hasse invariant in Weierstraß form (A1.13)]
\label{prop:hasse-weierstrass}
\lean{AnalyticTheory.hasseInv_weierstrass}
\uses{def:hasse-invariant, def:normal-forms}
$\gamma\equiv-\tfrac12 g_2/g_3$ modulo squares.
\end{proposition}

\begin{lemma}[Lemma A (unique nontrivial automorphism)]
\label{lem:A}
\lean{AnalyticTheory.lemmaA}
\uses{def:absolute-invariant}
\notready
For $j\ne 0,12^3$ there is only one nontrivial automorphism fixing the prime of degree $1$. The fixing predicate is a placeholder.
\end{lemma}

\begin{remark}[Remark 1 (formula for $j$ from $y^2+xy=x^3+Bx+C$)]
\label{rem:A-1}
\lean{AnalyticTheory.remark_A1_absinv}
\uses{def:absolute-invariant}
$j=(1-48B)^3/\Delta$ in any characteristic.
\end{remark}

\begin{remark}[Remark 2 (separably closed $K$)]
\label{rem:A-2}
\lean{AnalyticTheory.remark_A2_sep_closed}
\uses{def:hasse-invariant}
Over a separably closed field the Hasse invariant is trivial.
\end{remark}

\begin{remark}[Remark 3 (constant extensions)]
\label{rem:A-3}
\lean{AnalyticTheory.remark_A3_absinv, AnalyticTheory.remark_A3_hasse}
\uses{def:absolute-invariant, def:hasse-invariant, def:constant-extension}
Absolute invariant is preserved by constant extensions; Hasse invariant maps to the image class.
\end{remark}

\begin{theorem}[Uniqueness of $F\,|\,K$]
\label{thm:A-uniqueness}
\lean{AnalyticTheory.absolute_hasse_determine}
\uses{rem:A-3, def:hasse-invariant}
Fields with equal $j$ and $\gamma$ are $K$-isomorphic.
\end{theorem}

\begin{theorem}[Existence of $F\,|\,K$ with given $j,\gamma$]
\label{thm:A-existence}
\lean{AnalyticTheory.absolute_hasse_exists}
\uses{def:absolute-invariant, def:hasse-invariant}
For $j\ne 0,12^3$ and $\gamma\neq 0$ there is $F$ with invariants $j,\gamma$.
\end{theorem}

\begin{definition}[Absolute invariant of an arbitrary elliptic field (placeholder)]
\label{def:absolute-invariant-gen}
\lean{AnalyticTheory.absInvGen}
\uses{def:splits-over, def:elliptic}
\notready
The absolute invariant $j\in K$ of $E\,|\,K$. Placeholder.
\end{definition}

\section{Appendix \S A2: arbitrary elliptic fields}

\begin{definition}[Splitting over $L$]
\label{def:splits-over}
\lean{AnalyticTheory.SplitsOver}
\uses{def:constant-extension, def:degree-one-prime}
$E$ splits over $L$ if some constant extension $EL\,|\,L$ has a prime divisor of degree $1$.
\end{definition}

\begin{definition}[Hasse invariant of an arbitrary elliptic field (placeholder)]
\label{def:hasse-general}
\lean{AnalyticTheory.hasseGen}
\uses{def:splits-over, def:elliptic}
\notready
The Hasse invariant of $E\,|\,K$. Placeholder.
\end{definition}

\begin{theorem}[Independence of the splitting field]
\label{thm:A2-absinv}
\lean{AnalyticTheory.absInvGen_independent}
\uses{def:absolute-invariant-gen}
$j_L\in K$ does not depend on the splitting field.
\end{theorem}

\begin{theorem}[Formula (51)]
\label{thm:A2-51}
\lean{AnalyticTheory.hasseGen_compatible}
\uses{def:hasse-general}
\notready
The Hasse invariant comes from a unique class $\gamma\in K^\times/K^{\times 2}$. Partially formalized ($\exists!$ with a trivial conjunct).
\end{theorem}

\begin{definition}[Jacobian field]
\label{def:jacobian-field}
\lean{AnalyticTheory.IsJacobianField}
\uses{def:absolute-invariant-gen, def:hasse-general, def:hasse-invariant}
The Jacobian field of $E$: a field with a prime of degree $1$ having the same absolute and Hasse invariants.
\end{definition}

\begin{theorem}[Formulas (52) and (53)]
\label{thm:A2-52-53}
\lean{AnalyticTheory.jacobian_characterization_52, AnalyticTheory.jacobian_constant_extensions}
\uses{def:jacobian-field}
\notready
Characterization (52) and constant extensions (53). Recorded with placeholders ($\mathrm{True}$).
\end{theorem}

\begin{definition}[Characteristic cocycle (placeholder)]
\label{def:char-cocycle}
\lean{AnalyticTheory.charCocycle}
\uses{def:jacobian-field, def:splits-over}
\notready
$\zeta_E:G\to\mathrm{Aut}(FL\,|\,L)$. Placeholder.
\end{definition}

\begin{theorem}[Cocycle with values in translations]
\label{thm:A2-cocycle-translation}
\lean{AnalyticTheory.jacobian_iff_cocycle_translation}
\uses{def:char-cocycle}
\notready
$F$ is the Jacobian of $E$ iff the characteristic cocycle takes values in translations. Recorded as $\mathrm{True}$.
\end{theorem}

\begin{proposition}[Classification by $H^1(G,T)$]
\label{prop:A2-classification}
\lean{AnalyticTheory.a2_classification, AnalyticTheory.cohomological_invariant_independent}
\uses{thm:A2-cocycle-translation, thm:IX}
\notready
Elliptic fields with invariants $j,\gamma$ split by $L$ correspond to pairs $\pm\zeta\in H^1(G,T)$. Recorded as $\mathrm{True}$.
\end{proposition}
